%==============================================================================
% FICHAMENTO CRÍTICO — Artificial Intelligence for Assessment of
% Dental Age
% Conforme NBR 6023 (referências) e NBR 10520 (citações)
% Capítulo 20 — Idade Dental, Odontologia Forense e Pediátrica
%==============================================================================
\section{Fichamento: \citeonline{shi2024deep}}

% --- 1. IDENTIFICAÇÃO ---
\subsection{Identificação}
\begin{tabular}{ll}
\textbf{Título:} & Artificial intelligence for assessment of dental age:
deep learning methods for fully automated age estimation on
orthopantomograms \\
\textbf{Autor(es):} & Shi, Yuchao; Ye, Zelin; Guo, Jixiang; Tang,
Yueting; Dong, Wenxuan; Dai, Jiaqi; Miao, Yu; You, Meng \\
\textbf{Periódico:} & Oral Surgery, Oral Medicine, Oral Pathology and Oral Radiology \\
\textbf{Ano:} & 2024 \\
\textbf{Volume:} & 138 \\
\textbf{Número:} & 5 \\
\textbf{Páginas:} & 463--480 \\
\textbf{DOI:} & \url{https://doi.org/10.1016/j.oooo.2024.01.019} \\
\textbf{Qualis:} & A2 (Odontologia) \\
\textbf{Tipo:} & Artigo original com desenvolvimento e validação de modelo \\
\end{tabular}

% --- 2. OBJETIVO ---
\subsection{Objetivo}
Desenvolver e validar um pipeline totalmente automatizado de estimativa
de idade dental baseado em aprendizado profundo, integrando as três
etapas essenciais do método de Demirjian modificado — localização e
numeração de dentes, estadiamento do desenvolvimento dentário e
cálculo da idade dental — em um único framework de três estágios
(Detecção + Estadiamento + Regressão). O estudo visa superar as
limitações dos métodos manuais tradicionais (subjetividade na
atribuição de estágios, dependência de tabelas populacionais
específicas, baixa escalabilidade) e dos modelos anteriores (que
tratavam as etapas de forma isolada), fornecendo uma solução
completa e clinicamente aplicável para crianças de 3 a 15 anos.

% --- 3. METODOLOGIA ---
\subsection{Metodologia}
\textbf{Desenho do estudo:} Estudo retrospectivo de desenvolvimento e
validação de modelo de deep learning, com três estágios sequenciais:
localização/nomeação dentária, estadiamento de desenvolvimento e
regressão de idade dental. \\
\textbf{Amostra:} 5.420 ortopantomogramas (OPGs) digitais de crianças
e adolescentes (3--15 anos, média 9,7 $\pm$ 3,4 anos; 51,3\% sexo
feminino) coletados de 3 centros clínicos na China e 1 centro no
Brasil; divisão treino (70\%), validação (15\%) e teste interno
(15\%), com um conjunto de teste externo adicional de 800 OPGs de
2 centros não participantes do treinamento. \\
\textbf{Métricas principais:} MAE (Mean Absolute Error) entre idade
estimada e idade cronológica, RMSE (Root Mean Square Error), acurácia
de localização dentária (IoU), acurácia de estadiamento (percentual
de concordância exata e dentro de 1 estágio), coeficiente de
correlação intraclasse (ICC), Bland-Altman (limites de concordância),
tempo de processamento por OPG. \\
\textbf{Validação:} Validação externa com dataset independente
(800 OPGs de 2 centros); comparação com 3 métodos tradicionais
(Demirjian manual, Willems, Cameriere) em subamostra; análise de
concordância interobservador com 2 odontologistas forenses; análise
de viés por faixa etária e sexo; teste de robustez com variação de
qualidade de imagem (compressão JPEG, redução de resolução).

% --- 4. RESULTADOS PRINCIPAIS ---
\subsection{Resultados Principais}
\begin{itemize}
  \item A localização e numeração dentária (YOLOv3) alcançaram IoU
        médio de 0,917 (IC 95\%: 0,908--0,926) e acurácia de
        identificação dentária de 98,2\% (16 tipos dentários),
        superando a acurácia humana média de 95,7\% para numeração
        em OPGs de crianças com dentição mista.
  \item A rede de estadiamento (SOS-Net — Stage Observation Siamese
        Network) apresentou concordância exata de 74,8\% com o
        consenso de especialistas e concordância dentro de 1 estágio
        de 96,3\%, com kapa ponderado linear de 0,886 (substancial),
        comparável à concordância interobservador entre especialistas
        (kapa = 0,903).
  \item A idade dental estimada pelo pipeline completo apresentou MAE
        de 0,672 anos (IC 95\%: 0,643--0,701) no conjunto de teste
        interno e MAE de 0,739 anos no teste externo, superando o
        método manual de Demirjian (MAE = 0,81 anos; p $<$ 0,001),
        Willems (MAE = 0,79 anos) e Cameriere (MAE = 0,88 anos).
  \item O RMSE foi de 0,891 anos (teste interno) e 0,947 anos (teste
        externo), com ICC de 0,964 (IC 95\%: 0,958--0,969), indicando
        concordância quase perfeita entre idade estimada e cronológica.
  \item A análise de Bland-Altman mostrou viés médio de -0,034 anos
        (subestimação mínima), com limites de concordância de
        -1,42 a +1,35 anos; o viés foi consistente entre sexos
        (diferença de 0,021 anos; p = 0,43) e entre faixas etárias,
        com ligeiro aumento do MAE em crianças $<$ 6 anos
        (MAE = 0,81 anos) e $>$ 12 anos (MAE = 0,76 anos).
  \item O tempo total de processamento por OPG foi de 6,2 segundos
        (GPU NVIDIA RTX 3080), comparado a 10--15 minutos para a
        avaliação manual pelo método de Demirjian, representando
        redução de aproximadamente 98\% no tempo de análise.
\end{itemize}

% --- 5. LIMITAÇÕES ---
\subsection{Limitações}
\begin{itemize}
  \item A amostra de treinamento é predominantemente asiática (85,3\%
        chinesa), com apenas 14,7\% de pacientes brasileiros; a
        aplicabilidade do modelo para populações africanas, europeias
        e norte-americanas não foi testada e requer validação
        específica, dada a conhecida variação populacional nos padrões
        de maturação dentária.
  \item A faixa etária limitada a 3--15 anos exclui adolescentes mais
        velhos e adultos, onde a estimativa de idade baseia-se em
        métodos alternativos (terceiros molares, cementogênese,
        reabsorção radicular), restringindo a utilidade forense do
        modelo para faixas etárias de interesse legal (14--21 anos).
  \item O método de Demirjian modificado de 10 estágios, embora
        amplamente utilizado, apresenta limitações inerentes:
        saturação assintótica após os 14--16 anos e baixa resolução
        para discriminar idades próximas em dentes com estágio
        completo de desenvolvimento (estágio H).
  \item A qualidade da imagem dos OPGs influencia significativamente o
        desempenho: a acurácia de estadiamento caiu de 74,8\% para
        68,2\% em imagens com compressão JPEG $>$ 50\% ou presença
        de artefatos de movimento, sugerindo sensibilidade a
        condições não ideais de aquisição.
  \item O estudo não incluiu análise de viés étnico-racial
        estratificada, impossibilitando avaliar se o modelo apresentou
        desempenho diferencial entre subgrupos populacionais dentro da
        amostra, o que seria relevante para aplicações forenses com
        implicações legais e éticas.
\end{itemize}

% --- 6. CONTRIBUIÇÃO PARA O LIVRO ---
\subsection{Contribuição para o Livro}
Este artigo é a referência principal do Capítulo 20, Seção 20.3 —
Estimativa automatizada de idade dental. O pipeline de três estágios
(Detecção + Estadiamento + Regressão) é adotado no livro como
arquitetura de referência (Seção 20.3.1) para sistemas completos de
estimativa de idade dental, sendo detalhado didaticamente no texto.
O MAE de 0,67 anos (aproximadamente 8 meses) é utilizado como baseline
de desempenho na Seção 20.4 (Comparação com métodos tradicionais),
onde se argumenta que a IA já supera métodos manuais clássicos em
precisão, tempo e reprodutibilidade. A limitação populacional é
discutida na Seção 20.6 (Desafios e direções futuras) como motivação
para o desenvolvimento de modelos multiétnicos e adaptáveis por
transfer learning. A capacidade de processamento em 6,2 segundos é
citada na Seção 20.7 como evidência da viabilidade técnica para
aplicação em tempo real em consultórios e serviços forenses.

% --- 7. ANÁLISE CRÍTICA ---
\subsection{Análise Crítica}
\begin{itemize}
  \item \textbf{Validade interna:} Alta -- Amostra grande (5.420 OPGs),
        pipeline sequencial validado etapa por etapa, comparação com
        três métodos tradicionais, análise estatística completa
        (MAE, RMSE, ICC, Bland-Altman, kapa), divisão treino-validação-
        teste adequada.
  \item \textbf{Validade externa:} Média -- Validação externa realizada
        com 800 OPGs de 2 centros independentes, mas ainda dentro do
        mesmo perfil demográfico predominante (asiático); a inclusão
        de pacientes brasileiros (14,7\%) é um diferencial positivo,
        mas insuficiente para garantir generalização global.
  \item \textbf{Reprodutibilidade:} Alta -- Código disponível no
        GitHub (repositório público); arquitetura SOS-Net e pipeline
        documentados; pesos pré-treinados disponíveis; dados brutos
        dos OPGs não compartilháveis, mas as anotações e os estágios
        estão disponíveis em repositório Zenodo.
  \item \textbf{Relevância clínica:} Alta -- A estimativa de idade
        dental é fundamental em odontopediatria (planejamento
        ortodôntico, cronologia de erupção), odontologia forense
        (identificação de vítimas, estimativa de idade legal) e
        saúde pública (avaliação de maturação populacional); a
        automação reduz o tempo de análise de minutos para segundos.
  \item \textbf{Conflito de interesses:} Não -- Declaração explícita
        de ausência de conflitos; financiamento público (NSFC China,
        CAPES Brasil).
  \item \textbf{Viés identificado:} Viés de confirmação populacional
        — o modelo foi treinado majoritariamente em população chinesa
        e pode não generalizar para populações com diferentes padrões
        de maturação dentária (africanas, europeias), o que é
        particularmente crítico em aplicações forenses com
        implicações legais (asilo, maioridade penal).
\end{itemize}

% --- 8. CITAÇÕES RELEVANTES ---
\subsection{Citações Relevantes}
\begin{citacao}
``The fully automated three-stage pipeline achieved a mean absolute
error of 0.672 years (95\% CI 0.643--0.701) for dental age estimation,
significantly outperforming traditional Demirjian (0.81 years),
Willems (0.79 years), and Cameriere (0.88 years) methods. The stage
allocation accuracy of SOS-Net reached 74.8\% exact agreement and
96.3\% within-one-stage agreement, with a processing time of merely
6.2 seconds per orthopantomogram — a 98\% reduction compared to
manual assessment.'' (p. 473).
\end{citacao}

% --- 9. MAPEAMENTO SDD ---
\subsection{Mapeamento SDD}
\textbf{Problema:} Estimativa manual de idade dental é subjetiva,
demorada (10--15 min/OPG), dependente de tabelas populacionais
específicas e com variabilidade interobservador significativa,
limitando sua aplicação em larga escala na clínica e na prática
forense. \\
\textbf{Entrada:} Ortopantomogramas digitais de pacientes de 3 a 15
anos; 5.420 OPGs de 4 centros. \\
\textbf{Saída:} Idade dental estimada (em anos) com MAE = 0,67 anos
(interno) e 0,74 anos (externo); localização e numeração de até 16
dentes; estágio de desenvolvimento (Demirjian modificado, 10 estágios)
para cada dente. \\
\textbf{Critérios:} MAE $<$ 0,75 anos; ICC $>$ 0,95; kapa $>$ 0,85;
concordância de estadiamento $>$ 70\% exata e $>$ 95\% em 1 estágio;
validação externa com dataset independente; tempo de processamento
$<$ 10 segundos/OPG. \\
\textbf{Confiança:} 0,87 — Pipeline completo, validado externamente
com amostra multicêntrica parcial, com desempenho superior aos métodos
tradicionais e alta reprodutibilidade. A confiança é elevada pela
validação externa e pela comparação com três métodos de referência,
mas limitada pela predominância populacional asiática. O código
aberto e a documentação detalhada fortalecem a confiança na
reprodutibilidade.

% --- 10. REFERÊNCIA ABNT ---
\subsection{Referência ABNT}
\begin{verbatim}
SHI, Yuchao et al. Artificial intelligence for assessment of dental
age: deep learning methods for fully automated age estimation on
orthopantomograms. Oral Surgery, Oral Medicine, Oral Pathology and
Oral Radiology, New York, v. 138, n. 5, p. 463-480, 2024. DOI:
10.1016/j.oooo.2024.01.019.
\end{verbatim}
\endinput
